
%==============================================================
% PROJECT 15 — PILOTNET END-TO-END STEERING
% 8-PAGE AUTONOMY PORTFOLIO SECTION
%==============================================================
%
% REQUIRED IMAGE FILES — EACH USED EXACTLY ONCE
% -------------------------------------------------------------
% rawHistogram.png
% mirroredHistogram.png
% balancedHistogram.png
% TrainingANDvalidation.png
% EpochValLoss1.png
% EpochValLoss2.png
% EpochValLoss3.png
% ModelPredictionsPlot.png
%
% Assumes the portfolio already defines:
% DeepGreen, PolyGold, SoftGreen, SoftGold, Muted, Ink, RuleGray
% and loads TikZ, graphicx, hyperref, fontawesome5.
%
% Instructor objective reflected in this layout:
% • Predict a CONTINUOUS steering angle directly from camera imagery.
% • Analyze steering-angle imbalance before training.
% • Rebalance the dataset through binning / sampling.
% • Compare learning rate, batch size, and training duration.
% • Read validation-loss histories for stability and convergence.
%==============================================================

%--------------------------------------------------------------
% PAGE BASE
%--------------------------------------------------------------
\newcommand{\PNPageBase}{%
    \fill[white] (current page.south west) rectangle (current page.north east);
    \fill[DeepGreen]
      (current page.north west)
      rectangle
      ($(current page.north east)+(0,-0.095in)$);
    \fill[PolyGold]
      ($(current page.north west)+(0,-0.095in)$)
      rectangle
      ($(current page.north east)+(0,-0.14in)$);

    \fill[DeepGreen]
      (current page.south west)
      rectangle
      ($(current page.south east)+(0,0.36in)$);
    \fill[PolyGold]
      ($(current page.south west)+(0,0.36in)$)
      rectangle
      ($(current page.south east)+(0,0.405in)$);
}

%--------------------------------------------------------------
% HEADER
%--------------------------------------------------------------
\newcommand{\PNHeader}[3]{%
    \node[
        anchor=north west,
        text=PolyGold,
        font=\sffamily\bfseries\fontsize{8.1}{8.7}\selectfont
    ]
    at ($(current page.north west)+(0.42in,-0.39in)$)
    {03 AUTONOMY
     \hspace{5pt}/\hspace{5pt}
     PROJECT 15
     \ifx&#1&\else
     \hspace{5pt}/\hspace{5pt}
     #1
     \fi};

    \node[
        anchor=north west,
        text=DeepGreen,
        font=\bfseries\fontsize{21.0}{22.0}\selectfont
    ]
    at ($(current page.north west)+(0.40in,-0.67in)$)
    {#2};

    \node[
        anchor=north west,
        text=Muted,
        font=\sffamily\fontsize{8.15}{8.95}\selectfont
    ]
    at ($(current page.north west)+(0.42in,-1.12in)$)
    {#3};

    \fill[PolyGold]
      ($(current page.north west)+(0.42in,-1.36in)$)
      rectangle
      ($(current page.north east)+(-0.42in,-1.315in)$);
}

%--------------------------------------------------------------
% FOOTER
%--------------------------------------------------------------
\newcommand{\PNFooter}[3]{%
    \node[
        anchor=south west,
        text=white,
        font=\sffamily\bfseries\fontsize{7.0}{7.6}\selectfont
    ]
    at ($(current page.south west)+(0.42in,0.17in)$)
    {#1};

    \node[
        anchor=south,
        text=white,
        font=\sffamily\fontsize{7.0}{7.6}\selectfont
    ]
    at ($(current page.south)+(0,0.17in)$)
    {15 PILOTNET END-TO-END STEERING
     \hspace{6pt}
     #2};

    \node[
        anchor=south east,
        text=white,
        font=\sffamily\bfseries\fontsize{7.0}{7.6}\selectfont
    ]
    at ($(current page.south east)+(-0.42in,0.17in)$)
    {#3};
}

%--------------------------------------------------------------
% METRIC CARD
%--------------------------------------------------------------
\newcommand{\PNMetric}[6]{%
    \fill[SoftGreen,rounded corners=3pt]
      ($(current page.north west)+(#1,#2)$)
      rectangle ++(#3,-0.78in);

    \node[
        anchor=north west,
        text=DeepGreen,
        font=\sffamily\bfseries\fontsize{6.8}{7.4}\selectfont
    ]
    at ($(current page.north west)+(#1+0.16in,#2-0.12in)$)
    {#4};

    \node[
        anchor=north west,
        text=Ink,
        font=\bfseries\fontsize{13.0}{13.5}\selectfont
    ]
    at ($(current page.north west)+(#1+0.16in,#2-0.36in)$)
    {#5};

    \node[
        anchor=north west,
        text=Muted,
        font=\sffamily\fontsize{5.9}{6.5}\selectfont
    ]
    at ($(current page.north west)+(#1+0.16in,#2-0.61in)$)
    {#6};
}

%==============================================================
% PAGE 1 — PROJECT OBJECTIVE + RAW DATA PROBLEM
%==============================================================
\clearpage
\thispagestyle{empty}
\hypertarget{pilotnet}{}
\null

\begin{tikzpicture}[remember picture,overlay]
\PNPageBase

\PNHeader
{}
{PILOTNET END-TO-END STEERING}
{Cal Poly Autonomous Vehicles Laboratory
\hspace{7pt}\textbullet\hspace{7pt}
Three-Person Team
\hspace{7pt}\textbullet\hspace{7pt}
Spring 2026
\hspace{7pt}\textbullet\hspace{7pt}
San Luis Obispo, California}

% Intro
\node[
    anchor=north west,
    align=left,
    text width=10.00in,
    text=Ink,
    font=\fontsize{8.2}{9.5}\selectfont
]
at ($(current page.north west)+(0.48in,-1.62in)$)
{
The final autonomy project moved from explicit lane-state logic to
\textbf{end-to-end learning}. A PilotNet-style convolutional neural network
was trained to predict a \textbf{continuous steering angle directly from a
camera image}. Unlike classification, the output is a real-valued command,
so the distribution of steering examples strongly affects what behavior the
network learns.
};

% Main image
\node[
    anchor=north west,
    inner sep=1.4pt,
    draw=DeepGreen,
    line width=0.65pt
]
at ($(current page.north west)+(0.48in,-2.28in)$)
{
\includegraphics[
    width=6.18in,
    height=4.65in,
    keepaspectratio
]{\detokenize{rawHistogram.png}}
};

\fill[DeepGreen,rounded corners=1.5pt]
($(current page.north west)+(0.65in,-2.53in)$)
rectangle
($(current page.north west)+(2.60in,-2.28in)$);

\node[
    anchor=west,
    text=white,
    font=\sffamily\bfseries\fontsize{6.2}{6.8}\selectfont
]
at ($(current page.north west)+(0.80in,-2.405in)$)
{RAW STEERING DISTRIBUTION};

% Right narrative
\node[
    anchor=north west,
    text=DeepGreen,
    font=\sffamily\bfseries\fontsize{8.4}{9.1}\selectfont
]
at ($(current page.north west)+(6.92in,-2.30in)$)
{WHY BALANCING COMES FIRST};

\node[
    anchor=north west,
    align=left,
    text width=3.47in,
    text=Ink,
    font=\fontsize{7.45}{8.7}\selectfont
]
at ($(current page.north west)+(6.92in,-2.67in)$)
{
The raw dataset is dominated by steering values near zero. That is exactly
what a vehicle produces when most recorded driving is straight or only mildly
curved.

\vspace{7pt}

If the network is trained on this distribution unchanged, minimizing average
error rewards the model for predicting near-straight steering too often.
Large steering angles are rarer, so their mistakes contribute less frequently
to the training objective.

\vspace{7pt}

The first engineering task was therefore not model tuning. It was
\textbf{reshaping the data so left, straight, and right steering examples were
represented more deliberately}.
};

% Bottom flow
\fill[SoftGreen,rounded corners=4pt]
($(current page.north west)+(0.48in,-7.07in)$)
rectangle
($(current page.north west)+(10.42in,-7.95in)$);

\node[
    anchor=west,
    text=DeepGreen,
    font=\sffamily\bfseries\fontsize{7.0}{7.7}\selectfont
]
at ($(current page.north west)+(0.78in,-7.35in)$)
{END-TO-END LEARNING PIPELINE};

\node[
    anchor=west,
    text=Ink,
    font=\sffamily\bfseries\fontsize{7.15}{7.9}\selectfont
]
at ($(current page.north west)+(0.78in,-7.67in)$)
{CAMERA IMAGE
\hspace{9pt}$\rightarrow$\hspace{9pt}
CNN FEATURE EXTRACTION
\hspace{9pt}$\rightarrow$\hspace{9pt}
CONTINUOUS STEERING REGRESSION
\hspace{9pt}$\rightarrow$\hspace{9pt}
STEERING ANGLE};

\PNFooter
{\hyperlink{unet}{\faArrowLeft\ PROJECT 14}}
{1 / 5}
{DATA BALANCING \faArrowRight}

\end{tikzpicture}

%==============================================================
% PAGE 2 — MIRROR + BALANCE
%==============================================================
\clearpage
\thispagestyle{empty}
\hypertarget{pilotnet-balance}{}
\null

\begin{tikzpicture}[remember picture,overlay]
\PNPageBase

\PNHeader
{DATA PREPARATION}
{FROM A BIASED LOG TO A TRAINABLE STEERING DATASET}
{Mirroring
\hspace{7pt}\textbullet\hspace{7pt}
Eight steering bins
\hspace{7pt}\textbullet\hspace{7pt}
Undersampling
\hspace{7pt}\textbullet\hspace{7pt}
Balanced regression targets}

\node[
anchor=north west,align=left,text width=10.00in,text=Ink,
font=\fontsize{8.15}{9.45}\selectfont
]
at ($(current page.north west)+(0.48in,-1.60in)$)
{
The balancing workflow deliberately changed the steering-angle distribution
\textbf{before} model training. Horizontal mirroring first created opposite-sign
steering examples, reducing left/right bias. The steering range was then divided
into eight bins and sampled to the same frequency so that hard-left, straight,
and hard-right commands contributed more evenly to the regression problem.
};

% MIRRORED HISTOGRAM — enlarged
\node[anchor=north west,inner sep=1.1pt,draw=DeepGreen,line width=0.60pt]
at ($(current page.north west)+(0.48in,-2.18in)$)
{\includegraphics[width=4.86in,height=3.55in,keepaspectratio]{\detokenize{mirroredHistogram.png}}};

\fill[DeepGreen,rounded corners=1.5pt]
($(current page.north west)+(0.62in,-2.39in)$)
rectangle
($(current page.north west)+(2.12in,-2.18in)$);
\node[anchor=west,text=white,font=\sffamily\bfseries\fontsize{5.9}{6.4}\selectfont]
at ($(current page.north west)+(0.77in,-2.285in)$)
{1 \hspace{5pt} MIRROR};

% BALANCED HISTOGRAM — enlarged
\node[anchor=north west,inner sep=1.1pt,draw=DeepGreen,line width=0.60pt]
at ($(current page.north west)+(5.56in,-2.18in)$)
{\includegraphics[width=4.86in,height=3.55in,keepaspectratio]{\detokenize{balancedHistogram.png}}};

\fill[DeepGreen,rounded corners=1.5pt]
($(current page.north west)+(5.70in,-2.39in)$)
rectangle
($(current page.north west)+(7.34in,-2.18in)$);
\node[anchor=west,text=white,font=\sffamily\bfseries\fontsize{5.9}{6.4}\selectfont]
at ($(current page.north west)+(5.85in,-2.285in)$)
{2 \hspace{5pt} BALANCE};

% LEFT EXPLANATION
\fill[SoftGreen,rounded corners=4pt]
($(current page.north west)+(0.48in,-5.84in)$)
rectangle
($(current page.north west)+(5.16in,-7.28in)$);
\node[anchor=north west,text=DeepGreen,font=\sffamily\bfseries\fontsize{7.8}{8.5}\selectfont]
at ($(current page.north west)+(0.76in,-6.08in)$)
{WHY MIRRORING WAS NOT ENOUGH};
\node[anchor=north west,align=left,text width=4.08in,text=Ink,font=\fontsize{7.05}{8.25}\selectfont]
at ($(current page.north west)+(0.76in,-6.46in)$)
{
Mirroring makes left and right steering more symmetric, but it still leaves
the dominant central peak near zero. A network trained on that distribution
would continue seeing far more near-straight examples than large-angle turns.
};

% RIGHT EXPLANATION
\fill[SoftGreen,rounded corners=4pt]
($(current page.north west)+(5.56in,-5.84in)$)
rectangle
($(current page.north west)+(10.42in,-7.28in)$);
\node[anchor=north west,text=DeepGreen,font=\sffamily\bfseries\fontsize{7.8}{8.5}\selectfont]
at ($(current page.north west)+(5.84in,-6.08in)$)
{BINNING MAKES THE TARGET DISTRIBUTION EXPLICIT};
\node[anchor=north west,align=left,text width=4.20in,text=Ink,font=\fontsize{7.05}{8.25}\selectfont]
at ($(current page.north west)+(5.84in,-6.46in)$)
{
Eight steering ranges were populated evenly by sampling down to the least
frequent bin. This intentionally trades dataset size for a more uniform
learning signal across the full steering-angle range.
};

% NEXT-STEP STRIP — lowered slightly from the previous revision; entire block moves together
\fill[SoftGold,rounded corners=4pt]
($(current page.north west)+(0.48in,-7.39in)$)
rectangle
($(current page.north west)+(10.42in,-8.01in)$);
\fill[PolyGold]
($(current page.north west)+(0.48in,-7.39in)$)
rectangle
($(current page.north west)+(0.58in,-8.01in)$);
\node[anchor=west,text=DeepGreen,font=\sffamily\bfseries\fontsize{7.0}{7.7}\selectfont]
at ($(current page.north west)+(0.80in,-7.69in)$)
{NEXT: preserve this balanced target distribution when creating the training and validation subsets.};

\PNFooter
{\faArrowLeft\ RAW DATA}
{2 / 8}
{TRAIN / VALIDATION SPLIT \faArrowRight}

\end{tikzpicture}

%==============================================================
% PAGE 3 — TRAIN / VALIDATION SPLIT
%==============================================================
\clearpage
\thispagestyle{empty}
\hypertarget{pilotnet-split}{}
\null

\begin{tikzpicture}[remember picture,overlay]
\PNPageBase

\PNHeader
{DATASET SPLIT}
{BALANCING THE TRAINING AND VALIDATION TARGETS}
{Balanced steering bins
\hspace{7pt}\textbullet\hspace{7pt}
Training subset
\hspace{7pt}\textbullet\hspace{7pt}
Validation subset
\hspace{7pt}\textbullet\hspace{7pt}
Generalization check}

\node[
anchor=north west,align=left,text width=10.00in,text=Ink,
font=\fontsize{8.2}{9.5}\selectfont
]
at ($(current page.north west)+(0.48in,-1.60in)$)
{
Balancing the full dataset was only useful if that structure survived the
train/validation split. Samples from each steering bin were therefore assigned
to both subsets so the network learned from a broad steering range while the
validation data independently tested how well that behavior generalized.
};

\node[anchor=north,inner sep=1.4pt,draw=DeepGreen,line width=0.70pt]
at ($(current page.north)+(0in,-2.08in)$)
{\includegraphics[width=8.66in,height=4.72in,keepaspectratio]{\detokenize{TrainingANDvalidation.png}}};

\fill[DeepGreen,rounded corners=1.5pt]
($(current page.north west)+(2.32in,-2.34in)$)
rectangle
($(current page.north west)+(4.68in,-2.08in)$);
\node[anchor=west,text=white,font=\sffamily\bfseries\fontsize{6.2}{6.8}\selectfont]
at ($(current page.north west)+(2.48in,-2.21in)$)
{BALANCED TRAIN / VALIDATION DISTRIBUTIONS};

\fill[SoftGreen,rounded corners=4pt]
($(current page.north west)+(0.48in,-6.98in)$)
rectangle
($(current page.north west)+(5.18in,-8.02in)$);
\node[anchor=north west,text=DeepGreen,font=\sffamily\bfseries\fontsize{7.5}{8.2}\selectfont]
at ($(current page.north west)+(0.76in,-7.22in)$)
{WHY BOTH SUBSETS MATTER};
\node[anchor=north west,align=left,text width=4.10in,text=Ink,font=\fontsize{6.85}{7.95}\selectfont]
at ($(current page.north west)+(0.76in,-7.56in)$)
{
The training subset drives weight updates. The validation subset is held out
from those updates and acts as an independent signal for tuning, early
stopping, and detecting poor generalization.
};

\fill[SoftGold,rounded corners=4pt]
($(current page.north west)+(5.40in,-6.98in)$)
rectangle
($(current page.north west)+(10.42in,-8.02in)$);
\fill[PolyGold]
($(current page.north west)+(5.40in,-6.98in)$)
rectangle
($(current page.north west)+(5.51in,-8.02in)$);
\node[anchor=north west,text=DeepGreen,font=\sffamily\bfseries\fontsize{7.5}{8.2}\selectfont]
at ($(current page.north west)+(5.72in,-7.22in)$)
{WHAT THIS ENABLES NEXT};
\node[anchor=north west,align=left,text width=4.20in,text=Ink,font=\fontsize{6.85}{7.95}\selectfont]
at ($(current page.north west)+(5.72in,-7.56in)$)
{
With target imbalance controlled, changes in validation loss can be interpreted
more meaningfully as effects of model training rather than simply the original
steering-angle distribution.
};

\PNFooter
{\faArrowLeft\ DATA BALANCING}
{3 / 8}
{TRAINING STRATEGY \faArrowRight}

\end{tikzpicture}

%==============================================================
% PAGE 4 — WHAT WAS ACTUALLY TUNED
%==============================================================
\clearpage
\thispagestyle{empty}
\hypertarget{pilotnet-training}{}
\null

\begin{tikzpicture}[remember picture,overlay]
\PNPageBase

\PNHeader
{TRAINING STRATEGY}
{TURNING DATA BALANCE INTO A CONTROL MODEL}
{Continuous regression
\hspace{7pt}\textbullet\hspace{7pt}
Validation loss
\hspace{7pt}\textbullet\hspace{7pt}
Early stopping
\hspace{7pt}\textbullet\hspace{7pt}
Hyperparameter search}

\node[
    anchor=north west,
    align=left,
    text width=10.00in,
    text=Ink,
    font=\fontsize{8.0}{9.3}\selectfont
]
at ($(current page.north west)+(0.48in,-1.62in)$)
{
Once the steering distribution was controlled, the remaining question was
how to train the network stably. The project varied \textbf{learning rate},
\textbf{batch size}, and effective \textbf{training duration}. Validation loss
was used as the main comparison signal because the goal was not simply to fit
the training data, but to produce steering predictions that generalized to
unseen frames.
};

% Pipeline
\node[
    anchor=north west,
    text=DeepGreen,
    font=\sffamily\bfseries\fontsize{8.2}{8.9}\selectfont
]
at ($(current page.north west)+(0.48in,-2.28in)$)
{MODEL DEVELOPMENT LOOP};

\foreach \x/\title/\sub in {
0.48/CAMERA FRAMES/balanced images,
2.52/PILOTNET CNN/feature extraction,
4.56/REGRESSION HEAD/continuous output,
6.60/VALIDATION LOSS/generalization check,
8.64/RETRAIN/update settings}
{
    \fill[white,rounded corners=3pt]
      ($(current page.north west)+(\x in,-2.72in)$)
      rectangle
      ($(current page.north west)+(\x in+1.76in,-3.46in)$);

    \draw[DeepGreen,line width=0.55pt,rounded corners=3pt]
      ($(current page.north west)+(\x in,-2.72in)$)
      rectangle
      ($(current page.north west)+(\x in+1.76in,-3.46in)$);

    \node[
        anchor=north,
        align=center,
        text=DeepGreen,
        font=\sffamily\bfseries\fontsize{6.4}{7.0}\selectfont
    ]
    at ($(current page.north west)+(\x in+0.88in,-2.94in)$)
    {\title};

    \node[
        anchor=north,
        align=center,
        text=Muted,
        font=\sffamily\fontsize{5.8}{6.4}\selectfont
    ]
    at ($(current page.north west)+(\x in+0.88in,-3.19in)$)
    {\sub};
}

\foreach \x in {2.36,4.40,6.44,8.48}
{
    \node[
        text=PolyGold,
        font=\bfseries\fontsize{13}{13}\selectfont
    ]
    at ($(current page.north west)+(\x in,-3.09in)$)
    {$\rightarrow$};
}

% Experiment cards
\node[
    anchor=north west,
    text=DeepGreen,
    font=\sffamily\bfseries\fontsize{8.2}{8.9}\selectfont
]
at ($(current page.north west)+(0.48in,-3.92in)$)
{THREE VARIABLES DROVE THE SEARCH};

\fill[SoftGreen,rounded corners=4pt]
($(current page.north west)+(0.48in,-4.30in)$)
rectangle
($(current page.north west)+(3.58in,-6.18in)$);

\node[
    anchor=north west,
    text=PolyGold,
    font=\sffamily\bfseries\fontsize{7.0}{7.7}\selectfont
]
at ($(current page.north west)+(0.74in,-4.56in)$)
{01};

\node[
    anchor=north west,
    text=DeepGreen,
    font=\sffamily\bfseries\fontsize{8.0}{8.7}\selectfont
]
at ($(current page.north west)+(1.13in,-4.53in)$)
{LEARNING RATE};

\node[
    anchor=north west,
    align=left,
    text width=2.54in,
    text=Ink,
    font=\fontsize{7.0}{8.2}\selectfont
]
at ($(current page.north west)+(0.74in,-4.94in)$)
{
Controls how large each optimization update is. Too large can stall or
oscillate; too small can make convergence unnecessarily slow.
};

\fill[SoftGreen,rounded corners=4pt]
($(current page.north west)+(3.74in,-4.30in)$)
rectangle
($(current page.north west)+(6.84in,-6.18in)$);

\node[
    anchor=north west,
    text=PolyGold,
    font=\sffamily\bfseries\fontsize{7.0}{7.7}\selectfont
]
at ($(current page.north west)+(4.00in,-4.56in)$)
{02};

\node[
    anchor=north west,
    text=DeepGreen,
    font=\sffamily\bfseries\fontsize{8.0}{8.7}\selectfont
]
at ($(current page.north west)+(4.39in,-4.53in)$)
{BATCH SIZE};

\node[
    anchor=north west,
    align=left,
    text width=2.54in,
    text=Ink,
    font=\fontsize{7.0}{8.2}\selectfont
]
at ($(current page.north west)+(4.00in,-4.94in)$)
{
Changes how many samples contribute to one gradient update. The report found
the most useful range was approximately \textbf{64--128}.
};

\fill[SoftGreen,rounded corners=4pt]
($(current page.north west)+(7.00in,-4.30in)$)
rectangle
($(current page.north west)+(10.42in,-6.18in)$);

\node[
    anchor=north west,
    text=PolyGold,
    font=\sffamily\bfseries\fontsize{7.0}{7.7}\selectfont
]
at ($(current page.north west)+(7.26in,-4.56in)$)
{03};

\node[
    anchor=north west,
    text=DeepGreen,
    font=\sffamily\bfseries\fontsize{8.0}{8.7}\selectfont
]
at ($(current page.north west)+(7.65in,-4.53in)$)
{TRAINING DURATION};

\node[
    anchor=north west,
    align=left,
    text width=2.78in,
    text=Ink,
    font=\fontsize{7.0}{8.2}\selectfont
]
at ($(current page.north west)+(7.26in,-4.94in)$)
{
Early stopping limited wasted training, but the strongest runs still needed
roughly \textbf{300 epochs} to approach their best validation loss.
};

% Best run summary
\fill[SoftGold,rounded corners=4pt]
($(current page.north west)+(0.48in,-6.50in)$)
rectangle
($(current page.north west)+(10.42in,-8.02in)$);

\fill[PolyGold]
($(current page.north west)+(0.48in,-6.50in)$)
rectangle
($(current page.north west)+(0.59in,-8.02in)$);

\node[
    anchor=north west,
    text=DeepGreen,
    font=\sffamily\bfseries\fontsize{8.0}{8.7}\selectfont
]
at ($(current page.north west)+(0.82in,-6.78in)$)
{BEST REPORTED CONFIGURATION};

\PNMetric{0.82in}{-7.18in}{2.05in}{EPOCHS}{300}{long enough to plateau}
\PNMetric{3.05in}{-7.18in}{2.05in}{LEARNING RATE}{$2\times10^{-4}$}{best reported setting}
\PNMetric{5.28in}{-7.18in}{2.05in}{BATCH SIZE}{128}{within stable range}
\PNMetric{7.51in}{-7.18in}{2.05in}{TEST LOSS}{0.0470}{best reported test loss}

\PNFooter
{\faArrowLeft\ TRAIN / VALIDATION}
{4 / 8}
{LOSS CURVES \faArrowRight}

\end{tikzpicture}

%==============================================================
% PAGE 5 — EPOCH COUNT INVESTIGATION
%==============================================================
\clearpage
\thispagestyle{empty}
\hypertarget{pilotnet-loss}{}
\null

\begin{tikzpicture}[remember picture,overlay]
\PNPageBase

\PNHeader
{MODEL EXPERIMENTS}
{HOW LONG DID THE NETWORK NEED TO TRAIN?}
{Epoch count
\hspace{7pt}\textbullet\hspace{7pt}
Validation loss
\hspace{7pt}\textbullet\hspace{7pt}
Early stopping
\hspace{7pt}\textbullet\hspace{7pt}
Convergence}

\node[
    anchor=north west,
    align=left,
    text width=10.00in,
    text=Ink,
    font=\fontsize{8.1}{9.4}\selectfont
]
at ($(current page.north west)+(0.48in,-1.62in)$)
{
The first training question was whether additional epochs were still producing
useful learning or merely pushing the model toward overfitting. Two otherwise
similar runs were compared at approximately 200 and 300 epochs. Their validation
loss histories reveal whether the model had actually reached a plateau.
};

% Large EpochValLoss1 plot
\node[
    anchor=north west,
    inner sep=1.4pt,
    draw=DeepGreen,
    line width=0.70pt
]
at ($(current page.north west)+(0.82in,-2.24in)$)
{
\includegraphics[
    width=9.28in,
    height=4.72in,
    keepaspectratio
]{\detokenize{EpochValLoss1.png}}
};

\fill[DeepGreen,rounded corners=1.5pt]
($(current page.north west)+(1.00in,-2.50in)$)
rectangle
($(current page.north west)+(3.42in,-2.24in)$);

\node[
    anchor=west,
    text=white,
    font=\sffamily\bfseries\fontsize{6.15}{6.7}\selectfont
]
at ($(current page.north west)+(1.16in,-2.37in)$)
{EPOCH COUNT \hspace{5pt} 200 vs. 300};

% Bottom left interpretation
\fill[SoftGreen,rounded corners=4pt]
($(current page.north west)+(0.48in,-7.08in)$)
rectangle
($(current page.north west)+(5.26in,-7.92in)$);

\node[
    anchor=north west,
    text=DeepGreen,
    font=\sffamily\bfseries\fontsize{7.6}{8.3}\selectfont
]
at ($(current page.north west)+(0.76in,-7.27in)$)
{MORE EPOCHS WERE STILL USEFUL};

\node[
    anchor=north west,
    align=left,
    text width=4.18in,
    text=Ink,
    font=\fontsize{6.9}{8.0}\selectfont
]
at ($(current page.north west)+(0.76in,-7.5in)$)
{
Both curves follow the same broad decay, but the longer run continues lowering
validation loss after the shorter run ends. In this case, additional training
was still extracting useful structure rather than simply memorizing the data.
};

% Bottom right takeaway
\fill[SoftGold,rounded corners=4pt]
($(current page.north west)+(5.48in,-7.08in)$)
rectangle
($(current page.north west)+(10.42in,-7.92in)$);

\fill[PolyGold]
($(current page.north west)+(5.48in,-7.08in)$)
rectangle
($(current page.north west)+(5.59in,-7.92in)$);

\node[
    anchor=north west,
    text=DeepGreen,
    font=\sffamily\bfseries\fontsize{7.6}{8.3}\selectfont
]
at ($(current page.north west)+(5.80in,-7.27in)$)
{TRAINING LESSON};

\node[
    anchor=north west,
    align=left,
    text width=4.16in,
    text=Ink,
    font=\fontsize{6.9}{8.0}\selectfont
]
at ($(current page.north west)+(5.80in,-7.5in)$)
{
Early stopping should be treated as a safeguard, not as proof that fewer epochs
are always better. The validation history must first show that improvement has
actually stalled.
};

\PNFooter
{\faArrowLeft\ TRAINING STRATEGY}
{5 / 8}
{LEARNING RATE \faArrowRight}

\end{tikzpicture}

%==============================================================
% PAGE 6 — LEARNING RATE INVESTIGATION
%==============================================================
\clearpage
\thispagestyle{empty}
\hypertarget{pilotnet-learning-rate}{}
\null

\begin{tikzpicture}[remember picture,overlay]
\PNPageBase

\PNHeader
{MODEL EXPERIMENTS}
{HOW FAST SHOULD THE OPTIMIZER LEARN?}
{Learning rate
\hspace{7pt}\textbullet\hspace{7pt}
Optimization speed
\hspace{7pt}\textbullet\hspace{7pt}
Validation loss
\hspace{7pt}\textbullet\hspace{7pt}
Stable convergence}

\node[
    anchor=north west,
    align=left,
    text width=10.00in,
    text=Ink,
    font=\fontsize{8.1}{9.4}\selectfont
]
at ($(current page.north west)+(0.48in,-1.62in)$)
{
The next experiment isolated \textbf{learning rate}, which controls the size of
each weight update during gradient descent. Two similar PilotNet runs were
compared at $1\times10^{-4}$ and $2\times10^{-4}$. The goal was to determine
whether the slower setting was genuinely more stable or simply delaying useful
convergence.
};

% Large EpochValLoss2 plot
\node[
    anchor=north west,
    inner sep=1.4pt,
    draw=DeepGreen,
    line width=0.70pt
]
at ($(current page.north west)+(0.82in,-2.24in)$)
{
\includegraphics[
    width=9.28in,
    height=4.72in,
    keepaspectratio
]{\detokenize{EpochValLoss2.png}}
};

\fill[DeepGreen,rounded corners=1.5pt]
($(current page.north west)+(1.00in,-2.50in)$)
rectangle
($(current page.north west)+(3.72in,-2.24in)$);

\node[
    anchor=west,
    text=white,
    font=\sffamily\bfseries\fontsize{6.15}{6.7}\selectfont
]
at ($(current page.north west)+(1.16in,-2.37in)$)
{LEARNING RATE \hspace{5pt} $1e{-4}$ vs. $2e{-4}$};

% Interpretation
\fill[SoftGreen,rounded corners=4pt]
($(current page.north west)+(0.48in,-7.08in)$)
rectangle
($(current page.north west)+(5.26in,-7.92in)$);

\node[
    anchor=north west,
    text=DeepGreen,
    font=\sffamily\bfseries\fontsize{7.6}{8.3}\selectfont
]
at ($(current page.north west)+(0.76in,-7.27in)$)
{FASTER CONVERGENCE WITHOUT INSTABILITY};

\node[
    anchor=north west,
    align=left,
    text width=4.18in,
    text=Ink,
    font=\fontsize{6.9}{8.0}\selectfont
]
at ($(current page.north west)+(0.76in,-7.5in)$)
{
Within the tested range, the $2e{-4}$ run reduced validation loss substantially
faster. The slower run remained at a higher loss for much of training instead of
showing a clear stability advantage.
};

\fill[SoftGold,rounded corners=4pt]
($(current page.north west)+(5.48in,-7.08in)$)
rectangle
($(current page.north west)+(10.42in,-7.92in)$);

\fill[PolyGold]
($(current page.north west)+(5.48in,-7.08in)$)
rectangle
($(current page.north west)+(5.59in,-7.92in)$);

\node[
    anchor=north west,
    text=DeepGreen,
    font=\sffamily\bfseries\fontsize{7.6}{8.3}\selectfont
]
at ($(current page.north west)+(5.80in,-7.27in)$)
{WHY THIS MATTERS};

\node[
    anchor=north west,
    align=left,
    text width=4.16in,
    text=Ink,
    font=\fontsize{6.9}{8.0}\selectfont
]
at ($(current page.north west)+(5.80in,-7.5in)$)
{
Learning rate changes how quickly the optimizer traverses the loss landscape.
Too large can overshoot useful minima, while too small can leave the model
under-trained within the available epoch budget.
};

\PNFooter
{\faArrowLeft\ EPOCH COUNT}
{6 / 8}
{INITIALIZATION \faArrowRight}

\end{tikzpicture}

%==============================================================
% PAGE 7 — RUN-TO-RUN VARIABILITY
%==============================================================
\clearpage
\thispagestyle{empty}
\hypertarget{pilotnet-initialization}{}
\null

\begin{tikzpicture}[remember picture,overlay]
\PNPageBase

\PNHeader
{MODEL EXPERIMENTS}
{WHY THE SAME NETWORK CAN TRAIN DIFFERENTLY TWICE}
{Random initialization
\hspace{7pt}\textbullet\hspace{7pt}
Stochastic training
\hspace{7pt}\textbullet\hspace{7pt}
Convergence variability
\hspace{7pt}\textbullet\hspace{7pt}
Repeatability}

\node[
    anchor=north west,
    align=left,
    text width=10.00in,
    text=Ink,
    font=\fontsize{8.1}{9.4}\selectfont
]
at ($(current page.north west)+(0.48in,-1.62in)$)
{
A final comparison used the same model structure and hyperparameter settings but
different training runs. The resulting validation curves diverged considerably.
This exposed an important machine-learning reality: random weight initialization
and stochastic optimization can change how quickly an otherwise identical
network reaches a useful solution.
};

% Large EpochValLoss3 plot
\node[
    anchor=north west,
    inner sep=1.4pt,
    draw=DeepGreen,
    line width=0.70pt
]
at ($(current page.north west)+(0.82in,-2.24in)$)
{
\includegraphics[
    width=9.28in,
    height=4.72in,
    keepaspectratio
]{\detokenize{EpochValLoss3.png}}
};

\fill[DeepGreen,rounded corners=1.5pt]
($(current page.north west)+(1.00in,-2.50in)$)
rectangle
($(current page.north west)+(3.86in,-2.24in)$);

\node[
    anchor=west,
    text=white,
    font=\sffamily\bfseries\fontsize{6.15}{6.7}\selectfont
]
at ($(current page.north west)+(1.16in,-2.37in)$)
{SAME MODEL \hspace{5pt} DIFFERENT INITIALIZATION};

% Interpretation
\fill[SoftGreen,rounded corners=4pt]
($(current page.north west)+(0.48in,-7.08in)$)
rectangle
($(current page.north west)+(5.26in,-7.92in)$);

\node[
    anchor=north west,
    text=DeepGreen,
    font=\sffamily\bfseries\fontsize{7.6}{8.3}\selectfont
]
at ($(current page.north west)+(0.76in,-7.27in)$)
{TRAINING IS STOCHASTIC};

\node[
    anchor=north west,
    align=left,
    text width=4.18in,
    text=Ink,
    font=\fontsize{6.9}{8.0}\selectfont
]
at ($(current page.north west)+(0.76in,-7.5in)$)
{
One run begins improving much earlier than the other even though the architecture
and nominal settings are the same. A single training history therefore should
not be treated as perfectly repeatable evidence.
};

\fill[SoftGold,rounded corners=4pt]
($(current page.north west)+(5.48in,-7.08in)$)
rectangle
($(current page.north west)+(10.42in,-7.92in)$);

\fill[PolyGold]
($(current page.north west)+(5.48in,-7.08in)$)
rectangle
($(current page.north west)+(5.59in,-7.92in)$);

\node[
    anchor=north west,
    text=DeepGreen,
    font=\sffamily\bfseries\fontsize{7.6}{8.3}\selectfont
]
at ($(current page.north west)+(5.80in,-7.27in)$)
{ENGINEERING IMPLICATION};

\node[
    anchor=north west,
    align=left,
    text width=4.16in,
    text=Ink,
    font=\fontsize{6.9}{8.0}\selectfont
]
at ($(current page.north west)+(5.80in,-7.5in)$)
{
Repeated runs, fixed random seeds, and consistent evaluation data improve
reproducibility. Hyperparameter conclusions are stronger when they survive
variation in initialization rather than depending on one favorable run.
};

\PNFooter
{\faArrowLeft\ LEARNING RATE}
{7 / 8}
{PREDICTION RESULT \faArrowRight}

\end{tikzpicture}

%==============================================================
% PAGE 8 — MODEL OUTPUT + FINAL ENGINEERING TAKEAWAYS
%==============================================================
\clearpage
\thispagestyle{empty}
\hypertarget{pilotnet-result}{}
\null

\begin{tikzpicture}[remember picture,overlay]
\PNPageBase

\PNHeader
{MODEL RESULT}
{FROM CAMERA FRAME TO CONTINUOUS STEERING}
{Predicted vs. ground-truth steering
\hspace{7pt}\textbullet\hspace{7pt}
Regression behavior
\hspace{7pt}\textbullet\hspace{7pt}
Final training lessons}

\node[
    anchor=north west,
    align=left,
    text width=10.00in,
    text=Ink,
    font=\fontsize{8.0}{9.3}\selectfont
]
at ($(current page.north west)+(0.48in,-1.62in)$)
{
The best model was evaluated by comparing its continuous steering output with
the recorded ground-truth signal. The network captures the broad sequence of
left, straight, and right steering behavior, while the prediction remains
smoother and more quantized than the recorded command in several regions.
That difference is exactly why validation loss alone is not enough: the final
signal still has to be interpreted as a vehicle-control command.
};

% Prediction plot
\node[
    anchor=north west,
    inner sep=1.2pt,
    draw=DeepGreen,
    line width=0.65pt
]
at ($(current page.north west)+(0.48in,-2.22in)$)
{
\includegraphics[
    width=9.94in,
    height=3.42in,
    keepaspectratio
]{\detokenize{ModelPredictionsPlot.png}}
};

\fill[DeepGreen,rounded corners=1.5pt]
($(current page.north west)+(0.65in,-2.46in)$)
rectangle
($(current page.north west)+(2.95in,-2.20in)$);

\node[
    anchor=west,
    text=white,
    font=\sffamily\bfseries\fontsize{6.0}{6.6}\selectfont
]
at ($(current page.north west)+(0.80in,-2.33in)$)
{BEST MODEL \hspace{5pt} PREDICTION VS. GROUND TRUTH};

% Metrics
\PNMetric{0.48in}{-5.82in}{2.20in}{EPOCHS}{300}{best reported model}
\PNMetric{2.86in}{-5.82in}{2.20in}{LEARNING RATE}{$2e{-4}$}{faster useful convergence}
\PNMetric{5.24in}{-5.82in}{2.20in}{BATCH SIZE}{128}{within stable range}
\PNMetric{7.62in}{-5.82in}{2.80in}{TEST LOSS}{0.0470}{lowest reported test loss}

% Final lessons
\fill[SoftGold,rounded corners=4pt]
($(current page.north west)+(0.48in,-6.70in)$)
rectangle
($(current page.north west)+(6.95in,-7.88in)$);

\fill[PolyGold]
($(current page.north west)+(0.48in,-6.70in)$)
rectangle
($(current page.north west)+(0.59in,-7.88in)$);

\node[
    anchor=north west,
    text=DeepGreen,
    font=\sffamily\bfseries\fontsize{7.8}{8.5}\selectfont
]
at ($(current page.north west)+(0.82in,-6.94in)$)
{ENGINEERING TAKEAWAYS};

\node[
    anchor=north west,
    align=left,
    text width=5.76in,
    text=Ink,
    font=\fontsize{6.65}{7.55}\selectfont
]
at ($(current page.north west)+(0.82in,-7.1in)$)
{
\textbf{Data distribution is part of the controller.}
A regression model only learns the steering behaviors represented in its
training examples.

\vspace{2.5pt}

\textbf{Validation history must be interpreted, not merely minimized.}
A lower loss is meaningful only if the resulting steering signal is also
physically useful.

\vspace{2.5pt}

\textbf{Training is stochastic.}
Repeated runs of the same architecture can converge differently because the
weights begin from different random initializations.
};

% Resources
\fill[SoftGreen!34,rounded corners=4pt]
($(current page.north west)+(7.18in,-6.70in)$)
rectangle
($(current page.north west)+(10.42in,-7.88in)$);

\draw[RuleGray,line width=0.45pt,rounded corners=4pt]
($(current page.north west)+(7.18in,-6.70in)$)
rectangle
($(current page.north west)+(10.42in,-7.88in)$);

\node[
    anchor=north west,
    text=DeepGreen,
    font=\sffamily\bfseries\fontsize{7.6}{8.3}\selectfont
]
at ($(current page.north west)+(7.46in,-6.94in)$)
{PROJECT RESOURCES};

\node[
    anchor=north west,
    align=left,
    text=DeepGreen,
    font=\sffamily\bfseries\fontsize{6.9}{7.7}\selectfont
]
at ($(current page.north west)+(7.46in,-7.31in)$)
{
\faFilePdf\ PROJECT REPORT

\vspace{7pt}

\faCode\ COLAB / GITHUB
};

\PNFooter
{\faArrowLeft\ INITIALIZATION}
{8 / 8}
{\hyperlink{overview}{PORTFOLIO OVERVIEW \faArrowRight}}

\end{tikzpicture}

%==============================================================
% END PROJECT 15
%==============================================================
